% optimization/alignment/00_context_and_problem.tex

\section*{Motivation}

\begin{remark}
In practical applications, alignments are not only reconstructed from
data, but also designed, modified, and improved.
\end{remark}

\begin{remark}
This leads naturally to optimization problems in which alignment
parameters are adjusted to satisfy geometric, physical, and regulatory
constraints, as already emphasized in railway alignment literature
\cite{Kufver2000RailwayAlignment,Klein1998Trassierung}.
\end{remark}

\section{General Optimization Problem}

\begin{definition}[Alignment optimization problem]
An alignment optimization problem consists of:
\begin{itemize}
\item a parameterized alignment model,
\item an objective function,
\item a set of constraints.
\end{itemize}
\end{definition}

\begin{definition}[Objective function]
An objective function is a mapping
\[
J : \mathcal{A} \to \R
\]
that assigns a cost or quality measure to an alignment $\mathcal{A}$.
\end{definition}

\begin{remark}
Typical objectives include:
\begin{itemize}
\item minimizing curvature variation,
\item minimizing deviation from measurement data,
\item minimizing construction cost,
\item maximizing comfort.
\end{itemize}
Such objective classes are consistent with both general numerical
optimization literature \cite{Nocedal2006NumericalOptimization}
and railway-specific optimization approaches
\cite{Kufver2000RailwayAlignment}.
\end{remark}

\section{Parameterization}

\begin{remark}
The choice of parameters is central to the formulation of the
optimization problem.
\end{remark}

\begin{definition}[Parameter vector]
An alignment is described by a parameter vector
\[
\mathbf{x} \in \R^n
\]
encoding quantities such as:
\begin{itemize}
\item element lengths,
\item curvature values,
\item transition parameters,
\item geometric constraints.
\end{itemize}
\end{definition}

\begin{remark}
The sparse alignment model provides a natural parameterization by
element-wise variables.
\end{remark}

\TODO{Explicit parameter sets for fixed and transition elements.}

\section{Constraints}

\begin{definition}[Constraint]
A constraint is a condition of the form
\[
g_i(\mathbf{x}) \leq 0,
\qquad
h_j(\mathbf{x}) = 0.
\]
\end{definition}

\begin{remark}
Constraints arise from:
\begin{itemize}
\item geometric continuity,
\item engineering design rules,
\item physical limitations,
\item regulatory requirements.
\end{itemize}
\end{remark}

\subsection{Geometric Constraints}

\begin{itemize}
\item continuity of position and direction,
\item continuity of curvature.
\end{itemize}

\subsection{Engineering Constraints}

\begin{itemize}
\item minimum radius,
\item maximum cant,
\item limits on curvature derivatives.
\end{itemize}

\begin{remark}
Such conditions reflect standard railway design practice and are closely
related to design rules discussed in engineering literature and
standards \cite{Klein1998Trassierung,EN13803_2017}.
\end{remark}

\OPEN{Formal representation of engineering rules in constraint form.}

\section{Relation to Numerical Reconstruction}

\begin{remark}
Evaluation of the objective function requires reconstruction of the
alignment geometry from the parameter vector.
\end{remark}

\begin{remark}
Thus, optimization is intrinsically linked to numerical reconstruction,
as discussed in the previous chapter. From a numerical perspective, this
places alignment optimization within the broader context of differential
equations and numerical approximation
\cite{Hairer1993ODE,Press2007NumericalRecipes}.
\end{remark}

\section{Sequential Quadratic Programming}

\begin{remark}
Sequential Quadratic Programming (SQP) is a powerful method for solving
nonlinear constrained optimization problems
\cite{Nocedal2006NumericalOptimization}.
\end{remark}

\begin{definition}[SQP iteration]
An SQP method approximates the original problem by a sequence of
quadratic subproblems.
\end{definition}

\begin{remark}
Each iteration solves a quadratic approximation of the objective
function subject to linearized constraints.
\end{remark}

\TODO{Explicit formulation of SQP for alignment parameters.}

\RESEARCH{Adaptation of SQP to sparse alignment structures.}

\section{Alignment Optimization as Control Problem}

\begin{remark}
The curvature function $\kappaf(s)$ may be interpreted as a control
variable in a continuous system.
\end{remark}

\begin{definition}[Control formulation]
The alignment problem can be formulated as an optimal control problem
with state variables $(\gamma,\theta)$ and control $\kappaf$.
\end{definition}

\begin{remark}
This viewpoint is compatible with the curvature-centred perspective on
railway alignment advocated by Kufver
\cite{Kufver1997MathematicalDescription,Kufver2000RailwayAlignment}.
\end{remark}

\RESEARCH{Connection between alignment optimization and optimal control theory.}

\section{Network-Level Optimization}

\begin{remark}
Real-world railway systems involve not only single alignments, but
networks with branching and merging structures.
\end{remark}

\RESEARCH{Extension of optimization from single alignments to networks.}

\OPEN{Handling of switches, crossings, and parallel tracks.}

\section{Robustness and Uncertainty}

\begin{remark}
Optimization must account for uncertainty in input data.
\end{remark}

\RESEARCH{Robust optimization under uncertain geometry.}

\OPEN{Trade-off between optimality and robustness.}

\section{Implementation Perspective}

\begin{remark}
Practical optimization systems must balance:
\begin{itemize}
\item numerical stability,
\item computational efficiency,
\item model expressiveness.
\end{itemize}
\end{remark}

\begin{remark}
This balance is central both in generic numerical optimization
\cite{Nocedal2006NumericalOptimization}
and in railway-specific optimization settings
\cite{Kufver2000RailwayAlignment}.
\end{remark}

\TODO{Integration of optimization into the modelling workflow.}

\section{Historical Context}

\begin{remark}
Classical systems such as AXTRAN have demonstrated the feasibility of
optimization-based alignment design.
\end{remark}

\begin{remark}
Modern approaches aim to extend these ideas using more flexible models
and improved numerical methods, in line with the development from
classical railway design methods toward explicit optimization
formulations \cite{Kufver2000RailwayAlignment}.
\end{remark}

\TODO{Reconstruction and generalization of AXTRAN-like methods.}

